\documentclass[conference]{IEEEtran}
\IEEEoverridecommandlockouts

\usepackage{amsmath,amssymb,amsfonts}
\usepackage{graphicx}
\usepackage{booktabs}
\usepackage{multirow}
\usepackage{xcolor}
\usepackage{tikz}
\usetikzlibrary{shapes.geometric, arrows.meta, positioning, fit, backgrounds, calc}
\usepackage{pgfplots}
\pgfplotsset{compat=1.18}
\usepackage[hidelinks]{hyperref}
\usepackage{url}
\usepackage{orcidlink}
\usepackage{float}

\begin{document}

\title{Trustworthy RAG: An Evaluation Agent for Detecting Misinformation and Knowledge Poisoning in Generative AI Systems}

\author{
\IEEEauthorblockN{
Balkrishna Giri,
Md Toufique Hasan\,\orcidlink{0009-0005-8907-9898},
Jussi Rasku\,\orcidlink{0000-0002-4401-8013},
Muhammad Waseem\,\orcidlink{0000-0001-7488-2577},
and Pekka Abrahamsson\,\orcidlink{0000-0002-4360-2226}
}
\IEEEauthorblockA{
Faculty of Information Technology and Communication Sciences, Tampere University\\
Tampere, Finland\\
\resizebox{0.95\textwidth}{!}{{\ttfamily e-mail: \{balkrishna.giri, mdtoufique.hasan, jussi.rasku, muhammad.waseem, pekka.abrahamsson\}@tuni.fi}}
}
}

\maketitle

\begin{abstract}
Retrieval-Augmented Generation (RAG) improves large language model (LLM) reliability by grounding responses in external knowledge. However, RAG systems usually trust the retrieved content. This creates a \emph{Security-Reliability Gap}: high semantic relevance does not guarantee factual truth. Adversaries can exploit this gap through \emph{knowledge poisoning}, where malicious documents are inserted into the corpus to cause targeted misinformation. We propose an \emph{Evaluation Agent} that acts as middleware in the RAG pipeline. It combines Natural Language Inference (NLI)-based factual verification, a five-signal poison detector with relevance-weighted aggregation, and a \emph{Trust Index} $T = 0.4\,F + 0.35\,C + 0.25\,(1-P)$ with a non-linear dampener for high-contamination contexts. On TruthfulQA with Llama~3.3~70B, the agent reaches 91\% accuracy and 100\% precision, improving by 7 percentage points over a naive always-trust baseline. It detects instruction injection with 100\% recall, while entity-swap and subtle in-place edits remain difficult to detect because they provide few surface signals. Across three LLMs, the results show that the Trust Index depends on the generator. Generation style matters more than model size: a 7B model reaches 87\% accuracy, while a 35B model reaches 69\% because its hedged outputs cause false positives. The Trust Index remains discriminative for all models, with ROC-AUC from 0.73 to 0.81, and per-LLM threshold calibration restores baseline-competitive accuracy. Results are stable across repeated runs ($90.6{\pm}0.5\%$ over five runs), and detection does not change with retrieval depth ($K{=}3 \equiv K{=}5$). In a software-engineering use case, a secure-coding assistant over OWASP/CWE guidance, the agent blocks injected unsafe advice with F1 92\% but misses subtle weakening. We release the proposed approach, attack generator, and experimental artifacts\footnote{\url{https://github.com/your-username/your-repository}}.
\end{abstract}

\begin{IEEEkeywords}
retrieval-augmented generation, knowledge poisoning, trustworthy AI, natural language inference, misinformation detection, LLM security
\end{IEEEkeywords}

\input{sections/01_introduction.tex}
\input{sections/02_background.tex}
\input{sections/03_approach.tex}
%==============================================================================
\section{Experimental Design}
\label{sec:experimental-design}
%==============================================================================

\textbf{Datasets.} We use two public benchmarks. \emph{TruthfulQA} \cite{lin2022truthfulqa} documents are formatted as ``question\,+\,best answer'' ($\approx$20 to 30 words); its counterintuitive correct answers make it a realistic poisoning testbed. \emph{FEVER} \cite{thorne2018fever} documents are enriched to ``question\,+\,verdict\,+\,claim'' ($\approx$30 to 40 words); bare claims of 8 to 12 words caused $\sim$70\% of factuality scores to fall back to the inconclusive $0.5$ baseline, motivating the enrichment. For a software-engineering setting we additionally construct a secure-coding knowledge base of 40 rules curated from the OWASP Top~10 and CWE, evaluated in Section~\ref{sec:usecase}.

\textbf{Protocol.} Each run has two parts: 50 clean queries (Part~A) and the same 50 queries against a corpus with 30\% of documents poisoned (Part~B), for 100 samples per run. We evaluate four rule-based poisoning strategies: \emph{contradiction} (appends a contradicting statement), \emph{instruction injection} (appends an override directive), \emph{entity swap} (in-place replacement of entities/numbers, leaving no appended artifacts), and \emph{subtle manipulation} (false qualifiers/hedging). A \emph{mixed} setting assigns strategies randomly across poisoned documents.

\textbf{Models and metrics.} Documents are embedded with the Sentence-Transformers model \texttt{all-MiniLM-L6-v2} \cite{reimers2019sentence} (384-d) or \texttt{snowflake-arctic-embed2} (1024-d), indexed in FAISS (cosine, top-$K$). Generation uses \texttt{llama3.3:70b} (primary) or \texttt{qwen3.5:35b}; NLI runs on CPU. We report accuracy, precision, recall, F1, and \emph{trust score separation} $\Delta = \bar{T}_{\text{clean}} - \bar{T}_{\text{poisoned}}$, a threshold-independent measure of discriminability. The comparison point is a \emph{naive always-trust} RAG that never flags; under 30\% poisoning its accuracy ($\approx$85\%) reflects class imbalance, not detection, so $\Delta$ and F1 are the informative metrics. We add 95\% confidence intervals: Wilson intervals for proportions and a percentile bootstrap ($B{=}20\,000$) for F1.

\textbf{Ground truth.} A sample is labeled poisoned by document-level index alignment between the clean and poisoned corpora. This fits the paired benchmark but differs from a retrieval-grounded label that checks whether a poisoned document actually enters the top-$K$ set.

\input{sections/05_results.tex}
\input{sections/06_discussion.tex}
\input{sections/07_conclusion.tex}
\input{sections/08_references.tex}

\end{document}
